% ATTEMPT_STATUS: unresolved
% RESEARCH_GENERATED: 2026-10-01; GPT-6 Astra Ultra; individual mathematical attempt
% TOP_PROBLEM: {"id": 95, "title": "Product-Free Sets in Finite Groups", "category_id": 4, "category": {"id": 4, "name": "algebra", "display_name": "Algebra", "description": "Group theory, ring theory, field theory, and algebraic structures.", "slug": "algebra", "order_index": 4, "created_at": "2026-07-31T15:26:25.671Z"}, "status": "open"}
% FIRST_POSTED: 2026-10-01
% Imported from: unsolvedmath_additional_500.tex; UnsolvedMath ID 95
% !TeX program = lualatex
% Compile twice: lualatex 95.tex
% Self-contained: no external bibliography, images, or \input files.
% Numeric section labels and visible numbers are original UnsolvedMath IDs.
\documentclass[11pt,a4paper]{article}
\usepackage[margin=24mm,headheight=14pt]{geometry}
\usepackage{fontspec}
\setmainfont{Latin Modern Roman}
\setsansfont{Latin Modern Sans}
\setmonofont{Latin Modern Mono}
\usepackage{amsmath,amssymb,mathtools}
\usepackage{unicode-math}
\setmathfont{Latin Modern Math}
\usepackage{microtype}
\usepackage{enumitem,xurl,needspace}
\usepackage[dvipsnames]{xcolor}
\usepackage{hyperref}
\hypersetup{unicode=true,colorlinks=true,linkcolor=MidnightBlue,urlcolor=MidnightBlue,
 pdftitle={UnsolvedMath Problem 95},pdfauthor={Source-annotated compilation},bookmarksnumbered=true}
\usepackage{bookmark,tocloft,titlesec,fancyhdr}
\setlength{\cftsecnumwidth}{4.2em}
\cftsetpnumwidth{2.5em}
\cftsetrmarg{4em}
\titleformat{\section}{\Large\bfseries\raggedright}{\thesection}{0.7em}{}
\pagestyle{fancy}\fancyhf{}
\fancyhead[L]{\small\sffamily UnsolvedMath research attempt}
\fancyhead[R]{\small Original catalogue IDs}
\fancyfoot[C]{\thepage}
\setlength{\parindent}{0pt}
\setlength{\parskip}{5pt plus 1pt}
\setlength{\emergencystretch}{3em}
\setcounter{tocdepth}{1}\setcounter{secnumdepth}{1}
\setlist[itemize]{leftmargin=3em,itemsep=4pt,topsep=4pt,parsep=0pt}
\allowdisplaybreaks
\newcommand{\N}{\mathbb{N}}
\newcommand{\Z}{\mathbb{Z}}
\newcommand{\Q}{\mathbb{Q}}
\newcommand{\R}{\mathbb{R}}
\newcommand{\C}{\mathbb{C}}
\newcommand{\F}{\mathbb{F}}
\newcommand{\A}{\mathbb{A}}
\newcommand{\PP}{\mathbb{P}}
\newcommand{\E}{\mathbb{E}}
\newcommand{\eps}{\varepsilon}
\newcommand{\dd}{\,\mathrm d}
\newcommand{\sref}[2]{\hyperlink{source:#1:#2}{[S#2]}}
\newif\ifshowcatalogue
\showcataloguetrue
\newcommand{\cataloguescope}[1]{\ifshowcatalogue
 \paragraph{Statement recorded in the supplied source.}
 #1\fi}
\begin{document}
% UnsolvedMath ID 95; source catalogue code GREEN-005

\setcounter{section}{94}

\section{Small Product-Free Subsets in Finite Groups}\label{95}

{\small\sffamily UnsolvedMath ID 95 · Algebra}

\subsection{Definitions and mathematical statement}

\paragraph{Shared notation.}Unless a section says otherwise, $\N=\{1,2,\ldots\}$,
$\Z$ is the ring of integers, $\Q,\R,\C$ are the rational, real, and complex
fields, $[N]=\{1,\ldots,\lfloor N\rfloor\}$, and $\log$ is the natural
logarithm. For real functions of a parameter tending to infinity,
$f\ll g$ and $f=O(g)$ mean $|f|\leq Cg$ eventually for some fixed $C$;
$f\gg g$ reverses this comparison; $f=o(g)$ means $f/g\to0$;
$f\sim g$ means $f/g\to1$; and $f\asymp g$ means both order bounds.
A subscript on $O$ or $\ll$ specifies allowed constant dependence.
The expression $n^{o(1)}$ in an upper bound means growth smaller than
$n^\epsilon$ for each fixed $\epsilon>0$. Local definitions govern any
exception, finite-field convention, or use of the same symbol for another
object. An asymptotic claim is not automatically an effective algorithm.



A subset $A$ of a finite group $G$ is product-free when $A\cap AA=\varnothing$, where $AA=\{ab:a,b\in A\}$ and equal factors are allowed. Put $p(G)=\max\{|A|:A\subseteq G,\ A\cap AA=\varnothing\}$. Extremal comparisons concern this maximum relative to the order of $G$, for example $\min_{|G|=N}p(G)$, and the groups attaining small values. Commutativity is not assumed. \sref{95}{1} \sref{95}{2}

\paragraph{Statement recorded in the supplied source.}
Which finite groups have the smallest largest product-free sets? \sref{95}{1}

\paragraph{Scope and interpretation.}

The source poses an extremal classification rather than a single numerical inequality. Fixing the group order makes the comparison meaningful; no new proposed asymptotic formula is inserted. \sref{95}{1}

\paragraph{Residual task reported by the archive.}

The embedded review (2026-08-17) records: Identify extremal groups and the true minimum order of the largest product-free subset. \sref{95}{1}

\subsection{Short English statement}

Which finite groups make even their largest product-free subsets unusually small?

\subsection{Research attempt}

\paragraph{Provenance and scope.}
This individual attempt was generated by GPT-6 Astra Ultra on 1 October 2026. The method was to derive explicit partial results and test the first proposed extension adversarially. The original statement and sources below are retained. No claim of mathematical novelty or complete solution is made.

\paragraph{Formal target and literature check.}
Write $p(G)$ for the maximum size of a product-free subset of the finite group $G$, and $\beta(G)=p(G)/|G|$. Equal factors are allowed, and the identity cannot belong to a product-free set. The author-hosted problem list confirms the extremal classification target. The following argument is elementary; it is not a replacement for the substantially stronger representation-theoretic bounds discussed there. We attack the quotient structure first, because a group with a large abelian quotient cannot be extremal in the strongest sense.

\paragraph{Quotient lifting, with the quantifiers checked.}
If $\pi:G\to Q$ is a surjective homomorphism and $B\subset Q$ is product-free, then $\pi^{-1}(B)$ is product-free. Indeed, $x,y,xy\in\pi^{-1}(B)$ would give $\pi(x),\pi(y),\pi(x)\pi(y)\in B$. Every fibre has size $|\ker\pi|$, so
\[
 p(G)\geq |\ker\pi|p(Q),\qquad \beta(G)\geq\beta(Q).
\]
This proves a necessary condition for a sequence with $\beta(G_i)\to0$: its nontrivial quotients cannot come from any fixed finite collection. More quantitatively, if $Q$ has order at most $M$ and is nontrivial, choose any $q\ne1$. The singleton $\{q\}$ is product-free, since $q^2=q$ would imply $q=1$. Hence $\beta(G)\geq1/M$. This argument uses a quotient, not merely a subgroup; replacing the quotient by a subgroup loses the fibre factor and does not give the stated density bound.

\paragraph{The half-density boundary can be classified exactly.}
For any nonempty product-free $A$ and any $a\in A$, the sets $A$ and $aA$ are disjoint and have equal cardinality. Therefore $|A|\leq |G|/2$. Suppose equality holds, and put $B=G\setminus A$. Then $aA=B$ for every $a\in A$. Fix $a_0\in A$. For every $a\in A$ we have $a_0^{-1}aA=A$. The left stabiliser
\[
 H=\{h\in G:hA=A\}
\]
is a subgroup and contains the $|A|$ distinct elements $a_0^{-1}a$. It is proper, since a nonempty proper subset cannot be invariant under all left translations: if $u\in A$, then every $v\in G$ equals $(vu^{-1})u$. Thus $|H|\geq |G|/2$ and $H\ne G$ force $[G:H]=2$. Also $A=a_0H$, since $a_0^{-1}A\subset H$ and both sets have size $|H|$. Finally $a_0\notin H$: otherwise $A=H$ would contain the identity. Conversely the nonidentity coset of any index-two subgroup is product-free, because its square is the subgroup. We have proved
\[
 p(G)=|G|/2\quad\Longleftrightarrow\quad G\text{ has an index-two subgroup}.
\]
Normality was not assumed: every index-two subgroup is normal, so the converse multiplication is justified.

\paragraph{Testing the tempting simple-group reduction.}
The quotient lemma suggests replacing $G$ by a simple quotient. Every nontrivial finite group has a maximal proper normal subgroup, so such a quotient exists. However its order can be much smaller than $|G|$. A bound $p(Q)\geq c|Q|^\theta$ with $\theta<1$ lifts to
\[
 p(G)\geq c|G||Q|^{\theta-1}\geq c|G|^\theta.
\]
The last inequality has the correct direction because the exponent is negative. Consequently a universal power lower bound for all finite simple groups would indeed imply the same lower bound for all finite groups. This is a useful rigorous reduction. It does not identify extremisers: a quotient substantially smaller than $G$ gives a stronger lower bound, while the actual maximal product-free set might be larger still. Nor does it supply a matching upper bound for any candidate simple group.

\paragraph{Verification and remaining gap.}
The companion exact script enumerates every subset of cyclic groups of orders two through nine, $S_3$, and $A_4$, checking every ordered product including squares. These are finite consistency checks for the definitions and half-density criterion, not evidence for an asymptotic exponent. The singleton argument also handles groups of prime order; the trivial group has $p(G)=0$ and is deliberately excluded from the simple-quotient reduction. The strongest result proved here is the exact half-density classification together with the transfer of any simple-group power bound. The unresolved step is a sharp bound on $p(S)$ for the relevant growing families of nonabelian finite simple groups, with upper bounds matching lower constructions. Nothing here establishes that one particular family minimises $p(G)$ at a given order.

\subsection{Sources}

{\small

\begin{itemize}

\item[\hypertarget{source:95:1}{S1}] ULAM.AI, \emph{UnsolvedMath}, supplied \texttt{problems.json}, record 95 (GREEN-005), statement and background. The embedded literature-review date is 2026-08-17; that is the archive’s review date, not a fresh certification. Dataset landing page: \url{https://huggingface.co/datasets/ulamai/UnsolvedMath}.

\item[\hypertarget{source:95:2}{S2}] Ben Green, 100 Open Problems (author’s maintained problem list). \url{https://people.maths.ox.ac.uk/greenbj/papers/open-problems.pdf}. The author-hosted document was inspected on 27 September 2026; the archive’s local code is not a problem-number locator in this paper.

\item[\hypertarget{source:95:3}{S3}] Ben Green problem collection, GREEN-005. \url{https://www.unsolvedmath.com/problems/GREEN-005}. Bibliographic information imported from the supplied record.

\end{itemize}

}

\label{end:95}
\end{document}
